\documentclass[11pt]{article}
\usepackage[margin=24mm]{geometry}
\usepackage{amsmath,amssymb,amsthm,mathtools,booktabs,hyperref}
\hypersetup{colorlinks=true,linkcolor=blue,urlcolor=blue,citecolor=blue}
\newtheorem{theorem}{Theorem}
\newtheorem{corollary}[theorem]{Corollary}
\newtheorem{proposition}[theorem]{Proposition}
\newcommand{\Herm}{\operatorname{Herm}}
\newcommand{\diag}{\operatorname{diag}}
\newcommand{\one}{\mathbf{1}}
\newcommand{\ii}{\mathrm{i}}
\newcommand{\status}[1]{\par\smallskip\noindent\textbf{Status: \texttt{#1}.}\ }
\title{Beyond Nodal Damping:\\Exact PLL retuning, collective modal spillover, and replacement limits}
\author{IEEE--39 research derivation}
\date{3 October 2026 --- theory first; no new dynamic simulations}
\begin{document}
\maketitle
\begin{abstract}
A finite change in one scalar PLL measurement delay produces a rank-one
change of the impedance at physical SG torque--speed ports. Its Hermitian
part has one positive and one negative eigenvalue if the reduced actuation
and measurement paths are not collinear. Thus latency is not generally a
positive-semidefinite damping penalty. Separately, a graph theorem relates
damping--inertia mismatch to a signed frequency-response moment through an
effective-resistance quadratic form. A negative moment forces overshoot in
a declared stable swing-network class. The implemented closed PLL transfer
also shows that pure delay first enters phase tracking at cubic order near
zero frequency. These results have explicit scopes: they are not a computed
replacement maximum, a nonlinear safety certificate, or a verified claim of
bibliographic novelty. The theory-only extension below derives a closed-form
PLL gain map that preserves a chosen oscillatory pole, an exact local law
for its effects on other modes, and a gain-bounded first-order replacement
limit. A further construction opens all PLLs simultaneously, eliminates
their gains analytically for one prescribed eigenpair, and gives exact
linear compatibility tests for several prescribed modal patterns.
It also identifies when permitted neighboring measurements can remove
that assignment conflict. These are design constructions, not evidence
of greater secure replacement. The fixed delays remain exogenous.
No new grid experiment is executed.
\end{abstract}

\section{Scientific contract and actual model}
The question is which physical interaction makes nodal damping insufficient
for SG$\to$GFL design, and what follows for stability and frequency security.
The focused extension in Section~\ref{sec:retuning} concerns exact
PLL retuning and its effects on other collective modes, followed by
a gain-bounded local replacement decision. The physical sign theorem
and graph moment below are supporting results with separate scopes.
Schur elimination, pole assignment, eigenvalue sensitivity and convex
separation are established tools; a derivation here is not evidence
of first publication. The graph moment's applicability to IEEE--39
remains to be established.

With fixed device support, write the nonlinear retarded DAE as
\begin{align}
 E(x,\rho)\dot x&=f_0(x,y;\rho)+\sum_i b_i(K_{p,i},K_{i,i})e_i(t-\tau_i),\\
 0&=g(x,y;\rho,d),\qquad
 e_i=-\sin\theta_i\,v_{r,i}+\cos\theta_i\,v_{i,i}. \label{eq:dae}
\end{align}
The exported reduced state convention has $E=I$. Physical inertia is
recovered by consistent rescaling of SG swing rows, not by mixing row
conventions. The implemented PLL is exactly
\begin{equation}
 \dot\theta_i=\omega_{p,i},\quad
 t_{f,i}\dot\omega_{p,i}=\xi_i+K_{p,i}e_i(t-\tau_i)-\omega_{p,i},
 \quad \dot\xi_i=K_{i,i}e_i(t-\tau_i). \label{eq:pll}
\end{equation}
Here $t_{f,i}=1/(300\,2\pi)$ s is the existing filter constant, distinct
from pure delay $\tau_i$. Other converter equations use $\theta_i$ but
do not directly use $\omega_{p,i}$ or $\xi_i$ in this implementation.
No current-loop, setpoint or DC signal is delayed by this contract.

Use initialized generator dispatch and rating:
\[
 \epsilon_i=1-\rho_i,\quad P_{SG,i}=\epsilon_iP_i^0,\quad
 P_{GFL,i}=\rho_iP_i^0,\quad M_i=2H_iS_i^0\epsilon_i,\quad
 J_{SG}=\sum_iP_i^0\epsilon_i .
\]
$S_i^0$ and $P_i^0$ are different quantities. Replacement changes
electrical injection/coupling as well as inertia. State removal at
$\rho_i=0$ or $1$ requires another support chart.

For consistent per-unit phasors $P_i+\ii Q_i=V_iI_i^*$.
The lossy network and voltage/current/DC dynamics remain inside $g,f$.
If $g_y$ is nonsingular, algebraic elimination gives
\[
 \delta y=-g_y^{-1}(g_x\delta x+g_d\delta d),\qquad
 A_{\rm inst}=f_x-f_yg_y^{-1}g_x
\]
in ODE coordinates, with the analogous substitution in each detector.
Thus P/Q and voltage effects enter the actual characteristic operator;
a lossless Laplacian does not replace the AC Jacobian.
No converter current limiter is implemented, so no limiter safety is claimed.

Sources inspected are
\path{experiments/nonlinear_codesign_20261001/ReducedDAE.jl},
\path{src/bnd_model_expN/PDExactDesignN.jl}, and
\path{experiments/graph_gsp_codesign_20261003/DelayedEvents.jl}.
The intended DC convention is physical supply in \path{ReducedDAE.jl};
the older helper by itself uses a different convention.

\section{First correction: derive the closed PLL before assigning damping}
\status{EXACT\_IDENTITY}
At a locked equilibrium $\theta_i^*=\arg V_i^*$, let
$V_i^\circ=|V_i^*|>0$ and let $\phi_i$ be terminal phase perturbation.
The linearized detector is $\delta e_i=V_i^\circ(\phi_i-\theta_i)$.
Eliminating the PLL frequency and integrator gives, for zero history,
\begin{equation}
 H_i(s)=\frac{\theta_i(s)}{\phi_i(s)}
 =\frac{V_i^\circ(K_{p,i}s+K_{i,i})e^{-s\tau_i}}
 {s^2(1+t_{f,i}s)+V_i^\circ(K_{p,i}s+K_{i,i})e^{-s\tau_i}}. \label{eq:H}
\end{equation}
For $K_{i,i}>0$, direct Taylor division proves
\begin{equation}
 H_i(s)=1-\frac{s^2}{V_i^\circ K_{i,i}}
 +\frac{K_{p,i}/K_{i,i}-t_{f,i}-\tau_i}{V_i^\circ K_{i,i}}s^3+O(s^4).
 \label{eq:series}
\end{equation}
Pure delay first appears at cubic order in this phase-tracking transfer,
not in its constant, linear or quadratic coefficients.
Any subsequent reduction using analytic factors and inverses nonsingular
at zero preserves the $O(s^3)$ delay difference. Dividing an angle
operator by $s$ shifts that difference to $O(s^2)$.
Singular elimination can change the order. The previously ill-conditioned
full hidden block at $s=0$ prevents an unqualified application of this
argument to the full reduced grid.

Consequently $L_\tau=\operatorname{Sym}(B_cK_iT_\tau C)$ is not established
for this implemented closed PLL. It may arise for another loop or
approximation, which would need its own derivation.
At finite frequency the inverse tracking defect is exactly
\begin{align}
 Q_i(s)&=H_i(s)^{-1}-1
 =\frac{s^2(1+t_{f,i}s)e^{s\tau_i}}{V_i^\circ(K_{p,i}s+K_{i,i})},\\
 |Q_i(\ii\Omega)|&=
 \frac{\Omega^2\sqrt{1+\Omega^2t_{f,i}^2}}
 {V_i^\circ\sqrt{K_{i,i}^2+\Omega^2K_{p,i}^2}},\\
 \arg Q_i(\ii\Omega)&=\pi+\arctan(\Omega t_{f,i})
 -\arctan(\Omega K_{p,i}/K_{i,i})+\Omega\tau_i\pmod{2\pi}.
\end{align}
Delay rotates this defect without changing its magnitude. This defect
is not itself physical damping or a Laplacian.

\section{Physical collective damping and its finite latency law}
\status{EXACT\_IDENTITY}
The exact characteristic operator of the full linearization is
\[
 \Delta(s)=sE-A_0-\sum_i b_ic_i^T e^{-s\tau_i}.
\]
The two delayed PLL injections share one column $b_i$, with entries
$K_{p,i}/t_{f,i}$ and $K_{i,i}$. $c_i^T$ differentiates the actual
detector after algebraic voltage elimination.

In exported $E=I$ coordinates, retain per-unit SG speed deviations $\nu$
and eliminate all other states $z$ wherever $\Delta_{zz}$ is invertible:
\begin{equation}
 Z(s)=M(\Delta_{\nu\nu}-\Delta_{\nu z}\Delta_{zz}^{-1}\Delta_{z\nu}),
 \qquad d=Z(s)\nu,\qquad D_H(\Omega)=\Herm Z(\ii\Omega). \label{eq:Z}
\end{equation}
$d_i$ is mechanical torque expressed as MW at nominal speed, injected
as $d_i/M_i$ in the swing equation. For physical harmonic phasors,
\[
 \langle d^T\nu\rangle=\tfrac12\widehat\nu^*D_H(\Omega)\widehat\nu .
\]
This is incremental mechanical supply, not total AC power or a nonlinear
Lyapunov function. A steady harmonic response needs to exist; an unstable
system may only admit a formal particular phasor solution.

\begin{theorem}[One delayed channel redistributes physical damping]
Fix $\rho$, equilibrium, gains, ports and $\Omega>0$.
Change only $\tau_i$ to $\tau_i+\delta\tau_i\geq0$.
Assume both hidden blocks are invertible. At $s=\ii\Omega$, evaluate
the following factors at the old design:
\begin{align*}
 \kappa_i&=e^{-s\tau_i}-e^{-s(\tau_i+\delta\tau_i)},&
 h_i&=c_{z,i}^T\Delta_{zz}^{-1}b_{z,i},\\
 u_i&=M(b_{\nu,i}-\Delta_{\nu z}\Delta_{zz}^{-1}b_{z,i}),&
 v_i^*&=c_{\nu,i}^T-c_{z,i}^T\Delta_{zz}^{-1}\Delta_{z\nu},\\
 \gamma_i&=\kappa_i/(1+\kappa_i h_i),&
 a_i&=\gamma_i u_i .
\end{align*}
Then
\begin{equation}
 \boxed{\delta Z=\gamma_i u_iv_i^*,\qquad
 \delta D_H=\tfrac12(a_iv_i^*+v_ia_i^*).} \label{eq:rank}
\end{equation}
If $\gamma_i\ne0$ and $u_i,v_i$ are independent, $\delta D_H$ has
exactly one positive and one negative eigenvalue, all others zero.
\end{theorem}
\begin{proof}
The full change is $\kappa_i b_ic_i^T$. Applying Sherman--Morrison
inside the Schur complement proves (\ref{eq:rank}). Nonsingularity of
both hidden blocks makes $1+\kappa_i h_i\ne0$ by the determinant lemma.
The range of $\delta D_H$ is contained in the span of $a_i,v_i$.
Its two possibly nonzero eigenvalues are
\[
 \ell_\pm=\frac{\Re(v_i^*a_i)\pm
 \sqrt{\|a_i\|^2\|v_i\|^2-[\Im(v_i^*a_i)]^2}}{2}.
\]
Their product is
\[
 \ell_+\ell_-=-\frac{|\gamma_i|^2}{4}
 \left(\|u_i\|^2\|v_i\|^2-|v_i^*u_i|^2\right)<0
\]
by strict Cauchy--Schwarz. This proves both signs and rank two.
\end{proof}
This is a finite-change statement for the full linearization, not
merely a first-order approximation. It says that a local latency
change increases harmonic dissipation in some collective speed
directions and decreases it in others, under the stated independence.
Hence its negative is not a positive-semidefinite delay penalty either.
Collinear paths can instead give a rank-one change of either sign or
zero; phase-periodic increments can make $\kappa_i=0$.
Multiple channel changes may combine to have another sign pattern.
This theorem alone does not determine stability or nonlinear safety.
IEEE--39 path independence has not been tested in this theory-only work.

Differentiating the finite formula gives
\begin{equation}
 Z_{\tau_i}(s)=s e^{-s\tau_i}u_iv_i^* . \label{eq:derivative}
\end{equation}
Away from the imaginary axis $v_i^*$ denotes the displayed analytic
row, not an additional conjugation of an analytic function.

\section{Graph meaning of the physical directions}
\status{EXACT\_IDENTITY}
Use a verified symmetric physical Laplacian only to define a basis:
$M^{-1/2}L_PM^{-1/2}=U\Lambda U^T$. The lossy model in $Z$ stays intact.
Under physical mass normalization,
\[
 \widehat{\delta D}=U^TM^{-1/2}\delta D_HM^{-1/2}U
 =\tfrac12(\widehat a\,\widehat v^*+\widehat v\,\widehat a^*),
 \quad \widehat a=U^TM^{-1/2}a,\quad \widehat v=U^TM^{-1/2}v .
\]
Thus $\widehat{\delta D}_{k\ell}
=\tfrac12(\widehat a_k\overline{\widehat v_\ell}
+\widehat v_k\overline{\widehat a_\ell})$.
Sylvester's law preserves the positive/negative counts under this
invertible congruence. Both network paths determine the modal effect;
the spatial Fourier spectrum of $\tau$ alone cannot supply its sign.
Low graph modes are not automatically the right control subspace.
The rank-two claim concerns one channel's increment, not the full
damping matrix.

\section{A graph law for inertia and aggregate frequency}
\subsection{An explicit reduced class}
Assume, separately from the full-model theorem,
\begin{equation}
 \nu=\dot q,\qquad Z_r(s)\nu=d,\qquad
 Z_r(s)=sM+\mathcal D(s)+L/s,\quad
 \mathcal D(s)=D_0+sD_1+O(s^2). \label{eq:reduced}
\end{equation}
$M$ is positive diagonal, $L=BWB^T$ is a connected undirected
Laplacian with $W>0$, and $\mathcal D$ is real analytic near zero.
For physical per-unit speed, $q=\delta/\omega_0$; consequently
$L=\omega_0 L_{\rm angle}$. $M,D_0,L$ then have units MW\,s, MW,
and MW/s. This scaling is essential.
The class (\ref{eq:reduced}) has not been derived for the full lossy
IEEE--39 PLL/DC model; singular hidden states at zero cannot be ignored.

Define
\[
 m=\one^TM\one,\quad d_0=\one^TD_0\one>0,\quad
 u_0=M^{1/2}\one/\sqrt m,\quad
 M^{-1/2}LM^{-1/2}=U\diag(0,\Lambda_\perp)U^T .
\]
Write $\widetilde D=U^TM^{-1/2}\mathcal D M^{-1/2}U$ and
$Z_\perp=sI+\widetilde D_{\perp\perp}+\Lambda_\perp/s$.
For inertia-proportional input $d=(M\one/m)p$, the COI speed
$w=\one^TM\nu/m$ satisfies
\begin{equation}
 z_{\rm coi}(s)=m[s+\widetilde D_{00}
 -\widetilde D_{0\perp}Z_\perp^{-1}\widetilde D_{\perp0}],
 \qquad w=p/z_{\rm coi}. \label{eq:coi}
\end{equation}
This is an exact Schur identity for this class wherever defined.

\begin{theorem}[Graph correction to the aggregate response moment]
Let $b_R=D_0\one-(d_0/m)M\one$ and
$b_L=D_0^T\one-(d_0/m)M\one$. Then
\begin{equation}
 z_{\rm coi}(s)=d_0+s m_{\rm app}+O(s^2),\qquad
 \boxed{m_{\rm app}=m+\one^TD_1\one-b_L^TL^\dagger b_R.} \label{eq:moment}
\end{equation}
If $D_0$ is symmetric, $b_L=b_R=b$ and
\begin{equation}
 \boxed{b^TL^\dagger b
 =m\sum_{k>0}\frac{\widetilde D_{k0}(0)^2}{\lambda_k}
 =\min_{Bf=b} f^TW^{-1}f\geq0.} \label{eq:flow}
\end{equation}
\end{theorem}
\status{PROVED\_REDUCED\_MODEL}
\begin{proof}
$Z_\perp^{-1}=s(\Lambda_\perp+s\widetilde D_{\perp\perp}+s^2I)^{-1}
=s\Lambda_\perp^{-1}+O(s^2)$. Substitute in (\ref{eq:coi}).
Both $b_L,b_R$ sum to zero. The mass-normalized spectral solution of
$Lx=b_R$ differs from $L^\dagger b_R$ only by a constant vector,
which $b_L^T$ annihilates. This identifies the cross coefficient as
$b_L^TL^\dagger b_R$. Symmetry gives the sum of squares.
For the flow identity, $f_*=WB^TL^\dagger b$ is feasible and
every feasible $f=f_*+c$ has $Bc=0$.
Its cost is $f_*^TW^{-1}f_*+c^TW^{-1}c$, proving the minimum.
\end{proof}
$m_{\rm app}$ is a response moment, not stored kinetic inertia.
It may be negative with positive physical inertia in a stable system.
The mismatch uses row/column sums of collective damping, not its
diagonal alone. Weak graph modes amplify the effect by $1/\lambda_k$.
Cycle flows belong to $\ker B$ and add the nonnegative cost above;
no separate cycle damping object is needed.
At fixed $b$, $\partial(b^TL^\dagger b)/\partial w_e
=-(B_e^TL^\dagger b)^2\leq0$.
This concerns the moment at a fixed operating point, not monotonic
improvement of nadir or stability under actual line changes.

\begin{corollary}[Negative moment forces an overshoot]
Suppose the frequency response is stable with integrable transient.
For $p(t)=-P1_{t\geq0}$, $P>0$, define signed frequency drop $x=-w$.
Then $x_\infty=P/d_0$ and
\begin{equation}
 \boxed{\int_0^\infty(x_\infty-x(t))\,dt=P m_{\rm app}/d_0^2.} \label{eq:area}
\end{equation}
If $m_{\rm app}<0$, the nadir is below the final frequency.
\end{corollary}
\begin{proof}
Expand $X(s)=P/[s z_{\rm coi}(s)]$.
The limit of $P/(d_0s)-X(s)$ at zero is the stated area.
A negative integral cannot have an everywhere nonnegative integrand.
\end{proof}
A positive moment does not rule out overshoot. A negative one does
not quantify the peak or establish violation of a declared limit.

\paragraph{Exact two-node example; no time simulation.}
Let $M=I_2$, $L=\left[\begin{smallmatrix}1&-1\\-1&1\end{smallmatrix}\right]$
in dimensionless units. Both $D_0=5I_2$ and $D_0=\diag(1,9)$ have
$m=2,d_0=10$. The former has $w/p=1/(2s+10)$, monotone.
For the latter, $b=(-4,4)^T$, $b^TL^\dagger b=16$, $m_{\rm app}=-14$,
and
\[
 \frac{w}{p}=\frac{s^2+5s+2}{2(s^3+10s^2+11s+10)}.
\]
The cubic is Hurwitz by $10,11,10>0$ and $10\cdot11>10$.
Equation (\ref{eq:area}) proves overshoot. Equal total inertia and
total damping do not determine the frequency transient.
This example is not an IEEE--39 result.

For a general deficit input $d(t)=-v1_{t\geq0}$ with $\one^Tv=P>0$,
put $a=v-(P/m)M\one$. Eliminating the disagreement input yields
\begin{equation}
 \int_0^\infty(x_\infty-x(t))\,dt
 =\frac{P m_{\rm app}}{d_0^2}+\frac{b_L^TL^\dagger a}{d_0}. \label{eq:location}
\end{equation}
The extra term comes from the numerator $P-s b_L^TL^\dagger a+O(s^2)$.
Thus the equations accept arbitrary input location. Algebraic load
events need their real input map, not an assumed mechanical step.

\section{Nonlinear causality, RoCoF and replacement}
\begin{proposition}[A causal limitation on PLL retuning]
Assume the nonlinear DDAE is well posed, its pre-event detector history
is zero, its trim is gain-independent, and gains occur only as in
(\ref{eq:pll}). At fixed other parameters, changing site $i$'s PLL
gains cannot change the physical trajectory for $0\leq t<\tau_i$.
If all tunable channels have positive delay, the trajectory before
$\tau_{\min}$ is independent of all their gains.
\end{proposition}
\status{EXACT\_IDENTITY}
\begin{proof}
The corresponding delayed error is zero on the stated interval.
Both designs have the same vector field and history there;
uniqueness of the retarded initial-value problem gives equality.
\end{proof}
Other converter states, currents, DC dynamics, SG governors and
algebraic voltages can respond immediately. This is not a claim
that GFL power is frozen before the PLL delay.
Any output-guard failure in that common interval is unrepairable
by these gains. If $\mathcal S_{\rm blind}$ is the set of $\rho$
passing all declared event/output guards on that interval, then
\[
 \{\rho:\exists K\ \text{fully feasible}\}\subseteq\mathcal S_{\rm blind},
 \qquad
 P_{GFL}^*\leq\sup_{\rho\in\mathcal S_{\rm blind}}\sum_iP_i^0\rho_i .
\]
This is an unevaluated necessary-condition relaxation, potentially
weak for short delays and long measurement windows.

For a mechanical step in (\ref{eq:reduced}), with causal controller
realization and no instantaneous synthetic-inertia feedthrough,
\[
 \dot\nu(0^+)=M^{-1}d(0^+),\qquad
 \dot w(0^+)=\one^Td(0^+)/m .
\]
An instantaneous COI rotor RoCoF limit $R_{\max}$ Hz/s therefore needs
\begin{equation}
 \sum_i2H_iS_i^0\epsilon_i
 \geq f_0|\one^Td(0^+)|/R_{\max}. \label{eq:floor}
\end{equation}
This classical condition is not novel. Minimizing $P^T\epsilon$
under it and $0\leq\epsilon\leq1$ is fractional knapsack: sort
$2H_iS_i^0/P_i^0$ decreasingly for positive dispatch and retain in
that order, with at most one final fraction. This gives only the
necessary-condition bound for that explicit input/output contract.

The frozen repository instead monitors bus-voltage phase over
$T_w=0.5$ s:
\[
 F_j(t)=\frac{\phi_j(t)-\phi_j(t-T_w)}{2\pi T_w},\quad
 R_j(t)=\frac{\phi_j(t)-2\phi_j(t-T_w)+\phi_j(t-2T_w)}{2\pi T_w^2}.
\]
Z-load events can jump algebraic voltage phase.
These metrics cannot be replaced by instantaneous rotor RoCoF.
No numerical upper replacement bound for the frozen problem is
computed in this document.

\section{The exact connection to stability and local design}
For a simple root, with nonzero eigenvalue denominator,
\[
 \lambda_p=-\frac{w^*\Delta_pv}{w^*\Delta_s v},\qquad
 \Delta_s=E+\sum_i\tau_i b_ic_i^T e^{-s\tau_i}.
\]
Delayed memory is not stored inertia. At a regular physical-port root,
\begin{equation}
 \boxed{\lambda_{\tau_i}
 =-\frac{\lambda e^{-\lambda\tau_i}
             (\ell^*u_i)(v_i^*r)}{\ell^*Z_s r}.} \label{eq:pole}
\end{equation}
This follows by using (\ref{eq:derivative}) in
$\lambda_p=-\ell^*Z_pr/(\ell^*Z_sr)$.
Phase and left/right mode overlaps matter; the harmonic sign theorem
does not replace root or nonlinear-security analysis.
The affine PLL injection also gives the consistency identity
$\lambda_{\tau_i}=-\lambda(\lambda_{\log K_{p,i}}+
\lambda_{\log K_{i,i}})$ at fixed $\rho$ and trim.
Multiple roots require subspace/nonsmooth analysis.

For $p=(\rho,\log K_p,\log K_i)$ a legitimate smooth local predictor is
\[
 \max_{\delta p}\sum_iP_i^0\delta\rho_i,\qquad
 g_j+\nabla_p g_j^T\delta p\leq0,
 \quad p+\delta p\text{ inside gain bounds and a trust region}.
\]
Spectral derivatives come from (\ref{eq:pole}); nonlinear peak
derivatives need the retarded variational equations and all active
times. This is standard sequential constrained optimization, not a
new global solver. No new optimized $\rho,K_p,K_i$ are produced here.

Nor should $m_{\rm app}\geq0$ be added as a necessary safety
constraint: overshoot below final frequency may still respect the
allowed frequency excursion. Positive harmonic damping at sampled
frequencies also does not prove stability.

\section{Exact retuning and its collective consequences}\label{sec:retuning}

\paragraph{Scope.}
The design question is whether additional replacement can be compensated
by bounded PLL gains, and which collective modes prevent that compensation.
We keep the actual nonlinear DDAE, its exact delayed linearization and
the physical ports. No new damping operator is postulated.
All delays remain externally specified. A change of delay below compares
two prescribed physical scenarios; it is not an optimization over latency.
Every formula uses a fixed state-support chart.

\subsection{The actual single-PLL return relation}
Remove the three PLL states at site $i$ and regard its angle $\theta_i$
as an input to the remaining grid. Keep all other PLLs, delays, SG,
voltage, current and DC equations. At a regular frequency write
\[
 \mathcal A_i(s;\rho)z=\mathcal B_i(s;\rho)\theta_i,\qquad
 e_i=\mathcal C_i(s;\rho)z+\mathcal D_i(s;\rho)\theta_i .
\]
The exact return relation is
\begin{equation}
 e_i=G_i(s;\rho)\theta_i,\qquad
 G_i=\mathcal D_i+\mathcal C_i\mathcal A_i^{-1}\mathcal B_i .
 \label{eq:return}
\end{equation}
The detector includes the negative local-angle term. $G_i$ is not
terminal voltage phase alone or a lossless scalar approximation.
It is independent of this site's removed $K_{p,i},K_{i,i},\tau_i$;
other controllers and delays are held fixed. Replacement dependence
includes the actual algebraic closure and equilibrium dependence.
Eliminating this PLL gives
\begin{equation}
 F_i(s)=s^2(1+t_{f,i}s)
 -(sK_{p,i}+K_{i,i})e^{-s\tau_i}G_i(s;\rho)=0 . \label{eq:scalarroot}
\end{equation}
Where $\mathcal A_i$ is nonsingular, the full determinant is
$\det(\mathcal A_i)F_i/t_{f,i}$ up to a nonzero ordering constant.
Hidden-block singularities must instead be treated with the original
characteristic operator. Not every pole need be visible in this relation.

\begin{theorem}[Exact single-site pole assignment with gain limits]
Fix $\rho$, all other controllers and delays, and
$\lambda=\alpha+\ii\omega$, $\omega>0$.
Assume $\mathcal A_i(\lambda)$ is invertible and
$G_i(\lambda;\rho)\ne0$. Define
\[
 W_i(\lambda,\rho,\tau_i)=
 \frac{\lambda^2(1+t_{f,i}\lambda)e^{\lambda\tau_i}}
 {G_i(\lambda;\rho)} .
\]
The unique real gains assigning $\lambda$ through this channel are
\begin{equation}
 \boxed{K_{p,i}=\frac{\Im W_i}{\omega},\qquad
 K_{i,i}=\Re W_i-\frac{\alpha}{\omega}\Im W_i .} \label{eq:gainexact}
\end{equation}
Under declared gain bounds, such an assignment is possible if and only if
\begin{equation}
 W_i\in\{\lambda k_p+k_i:\ k_p\in[K_p^-,K_p^+],\
                              k_i\in[K_i^-,K_i^+]\}. \label{eq:parallelogram}
\end{equation}
\end{theorem}
\status{EXACT\_IDENTITY}
\begin{proof}
Equation (\ref{eq:scalarroot}) reduces to
$\lambda K_{p,i}+K_{i,i}=W_i$.
Its imaginary and real parts give (\ref{eq:gainexact}), uniquely since
$\omega\ne0$. The image of the gain rectangle is precisely the
parallelogram in (\ref{eq:parallelogram}).
\end{proof}
This is a one-pole, one-site feasibility test with other gains frozen.
Failure is not a certificate against all distributed stabilization.
Pole assignment is established control theory; the system-specific
content is the exact return relation and the cross-mode result below.

\subsection{Transport between prescribed heterogeneous delays}
For frozen $\rho$ and an existing root $\lambda$, compare
$\tau_i$ with $\tau_i+h$. The exact gain map is
\begin{align}
 K_{p,i}(h)&=e^{\alpha h}\left[K_{p,i}\cos\omega h+
 \frac{K_{i,i}+\alpha K_{p,i}}{\omega}\sin\omega h\right],\\
 K_{i,i}(h)&=e^{\alpha h}\left[K_{i,i}\cos\omega h-
 \frac{|\lambda|^2K_{p,i}+\alpha K_{i,i}}{\omega}\sin\omega h\right].
 \label{eq:transport}
\end{align}
Indeed,
\[
 [\lambda K_{p,i}(h)+K_{i,i}(h)]e^{-\lambda(\tau_i+h)}
 =(\lambda K_{p,i}+K_{i,i})e^{-\lambda\tau_i}.
\]
The eliminated PLL equation at $\lambda$ is unchanged.
Applying this map to all prescribed heterogeneous delay changes
preserves the eliminated full-grid equations at
$\lambda,\overline\lambda$, for fixed $\rho$.
It does not ensure that these poles remain dominant or that the
other poles remain stable. All gain bounds must still hold.

The same map solves
\begin{equation}
 \frac{d}{dh}\begin{bmatrix}K_{p,i}\\K_{i,i}\end{bmatrix}
 =\begin{bmatrix}2\alpha&1\\-|\lambda|^2&0\end{bmatrix}
             \begin{bmatrix}K_{p,i}\\K_{i,i}\end{bmatrix}. \label{eq:gainflow}
\end{equation}
At $\alpha=0$, initially positive gains first lose positivity through
$K_i=0$ at $h=\omega^{-1}\arctan[K_i/(\omega K_p)]$ for increasing $h$.
Actual bounds may be reached earlier. This is a boundary of the
pole-preserving gain branch, not the maximum stable grid delay.

\begin{theorem}[Modal spillover under single-site pole preservation]
Let $\lambda=\alpha+\ii\omega$ and $\mu$ be simple characteristic
roots at the baseline, with return reductions regular at both.
Retune site $i$ by (\ref{eq:transport}) as its prescribed latency
changes, holding other sites fixed. Then
\begin{equation}
 \boxed{\left.\frac{d\mu}{dh}\right|_{0}
 =-K_{p,i}(\mu-\lambda)(\mu-\overline\lambda)
                         \frac{\partial\mu}{\partial K_{i,i}} .}
 \label{eq:spillover}
\end{equation}
The derivative on the right is in ordinary gain, not log gain.
\end{theorem}
\status{EXACT\_IDENTITY}
\begin{proof}
Differentiating (\ref{eq:scalarroot}) at a simple root gives
\[
 \mu_{K_{p,i}}=\mu\,\mu_{K_{i,i}},\qquad
 \mu_{\tau_i}=-\mu(K_{p,i}\mu_{K_{p,i}}+K_{i,i}\mu_{K_{i,i}}).
\]
Combining these with (\ref{eq:gainflow}),
\begin{align*}
 \dot\mu&=\mu_{\tau_i}
 +(K_{i,i}+2\alpha K_{p,i})\mu_{K_{p,i}}
 -|\lambda|^2K_{p,i}\mu_{K_{i,i}}\\
 &=-K_{p,i}(\mu^2-2\alpha\mu+|\lambda|^2)\mu_{K_{i,i}} .
\end{align*}
\end{proof}
Equivalently, for
$\mathcal K_i(s,h)=[sK_{p,i}(h)+K_{i,i}(h)]e^{-s(\tau_i+h)}$,
\[
 \partial_h\mathcal K_i(s,0)=
 -K_{p,i}e^{-s\tau_i}(s-\lambda)(s-\overline\lambda).
\]
For a finite change, $\mathcal K_i(s,h)-\mathcal K_i(s,0)$ has
zeros at $\lambda,\overline\lambda$, and is divisible by that
quadratic as an entire function of $s$.

\begin{corollary}[A precise single-site limitation]
If $\lambda_{K_{i,i}}\ne0$, any differentiable real retuning at
this site preserving the full complex root $\lambda$ while
$d\tau_i/dh=1$ has derivative (\ref{eq:gainflow}).
For $K_{p,i}>0$, it cannot also preserve a distinct root
$\mu\notin\{\lambda,\overline\lambda\}$ to first order when
$\mu_{K_{i,i}}\ne0$.
\end{corollary}
\begin{proof}
The gain-to-root real Jacobian is nonsingular:
$\lambda_{K_i}(\lambda\,\delta K_p+\delta K_i)$ has nonzero
real determinant for $\omega>0$, $\lambda_{K_i}\ne0$.
Compensation is therefore unique; (\ref{eq:spillover}) is nonzero.
\end{proof}
This is not an impossibility theorem for preserving only the real
parts of two roots, or for distributed retuning at other sites.
A complex root moving need not have a worsened real part.
The sign of $\Re\dot\mu$ is a prediction to test, not an assumed
destabilizing outcome.

\paragraph{Connection to Beyond Nodal Damping.}
The residue in (\ref{eq:spillover}) is
\[
 \mu_{K_{i,i}}=
 -\frac{\ell_\mu^*Z_{K_{i,i}}(\mu)r_\mu}{\ell_\mu^*Z_s(\mu)r_\mu}.
\]
The local PLL fixes the spectral polynomial; the lossy network,
inertia, P/Q and other devices fix its collective residue.
Nodal damping coefficients do not contain these left/right overlaps.
GSP can express the overlaps in a graph basis, but unsigned delay
roughness cannot substitute for their complex sign.

\subsection{Replacement at fixed physical delays}
At any prescribed $\rho$ in the regular region,
(\ref{eq:gainexact}) gives an exact candidate for one target pole.
For fixed $\lambda,\tau_i,t_{f,i}$,
\[
 W_{i,\rho_j}=-W_iG_{i,\rho_j}/G_i,\quad
 K_{p,i,\rho_j}=\Im W_{i,\rho_j}/\omega,\quad
 K_{i,i,\rho_j}=\Re W_{i,\rho_j}-\alpha\Im W_{i,\rho_j}/\omega .
\]
The derivative of $G_i$ must include equilibrium and algebraic-voltage
dependence. The sharing contract leaves $t_{f,i}$ fixed.
Holding a full complex pole is one candidate strategy.
For a requirement on real part alone, use
$\lambda=-\sigma_{\rm req}+\ii\omega$ with $\omega$ free;
fixing $\omega$ would add an unrequested design constraint.

\section{From modal prediction to a bounded replacement decision}

\begin{theorem}[Gain-bounded first-order replacement in a fixed direction]
Start at a feasible design. Fix $\rho=\rho_0+t r$, $t\ge0$,
with $P^Tr>0$, and an admissible interval $0\le t\le t_{\rm box}$.
Stack the included smooth constraint branches as $g\le0$ and put
$q=-g(p_0)\ge0$, $a=g_\rho r$, $C=g_K$.
Let the gain increments belong to a nonempty compact convex set
$\mathcal U$ containing zero. In the first-order constraint model,
\[
 \exists u\in\mathcal U:\quad t a+C u\le q
\]
if and only if
\begin{equation}
 t\,\eta^Ta\le\eta^Tq+h_{\mathcal U}(-C^T\eta)
                      \quad\text{for every }\eta\ge0 . \label{eq:dual}
\end{equation}
Consequently
\begin{equation}
 \boxed{t_{\max}^{\rm lin}
 =\min\left\{t_{\rm box},
 \inf_{\eta\ge0,\ \eta^Ta>0}
 \frac{\eta^Tq+h_{\mathcal U}(-C^T\eta)}{\eta^Ta}\right\}.}
 \label{eq:raybound}
\end{equation}
An empty infimum is $+\infty$. The maximal extra replacement
in this model is $t_{\max}^{\rm lin}P^Tr$ MW.
\end{theorem}
\status{PROVED\_REDUCED\_MODEL}
\begin{proof}
Feasibility means $q-ta\in C\mathcal U+\mathbb R_+^m$,
a closed convex set. Its separating halfspaces with finite
lower support use normals $\eta\ge0$, and
$\inf_{u\in\mathcal U}\eta^TCu=-h_{\mathcal U}(-C^T\eta)$.
This proves (\ref{eq:dual}). For $\eta^Ta\le0$ the inequality
is automatic because $q\ge0$, $0\in\mathcal U$ and $t\ge0$.
The remaining inequalities give (\ref{eq:raybound}).
\end{proof}
For the physical gain box $u_j\in[l_j,h_j]$,
$h_{\mathcal U}(z)=\sum_j\max(l_jz_j,h_jz_j)$.
The existing gain bounds are not enlarged.
If $C^T\eta=0$ for a nonnegative $\eta$, that weighted combination
of constraints has no first-order tuning authority, even if gain
increments were unbounded. Its replacement demand must be met
by baseline margin or another physical intervention.

This is standard convex separation applied to the problem; it
extends the repository's previous compensability formulation.
It is not a new general optimization theorem, a maximum over
all replacement directions, or a full nonlinear optimum.
The physical contribution sought is how the collective residues
determine the limiting combination of constraints.
For spectral maxima or trajectory peaks, retain all relevant
active branches or use directional nonsmooth analysis.
A finite list does not certify unseen roots, time instants or
disturbances.

For a uniformly verified remainder bound on the entire design region,
\[
 |g_j(p_0+(tr,u))-g_j(p_0)-t a_j-(Cu)_j|\le\varepsilon_j ,
\]
replacing $q$ by $q+\varepsilon$ gives a necessary relaxation,
and hence a regional upper bound on $t$.
Replacing $q$ by $q-\varepsilon$ in (\ref{eq:dual}) gives a
sufficient test for the included constraints. The simplified
ratio formula need not apply to $q-\varepsilon$, which can have
negative entries. No rigorous remainder bound has been obtained here.

\subsection{Analytic examples that expose the limiting mechanism}
Let $t\ge0$ denote normalized replacement and let $u\in[-\kappa,\kappa]$
be one normalized tuning direction. If the two modal constraints are
\[
 t+u\le q,\qquad t-u\le q,\qquad q\ge0,
\]
then adding them proves $t\le q$, attained with $u=0$.
Each constraint separately would allow $t=q+\kappa$.
The dual direction $\eta=(1,1)$ cancels the controller action and
detects the conflict. More tuning range does not remove this
particular first-order multimode limitation.
In contrast, if both rows are $t+u\le q$, the optimum is
$t=q+\kappa$, attained with $u=-\kappa$.
These exact algebraic examples explain what the test detects;
they do not assert that either row pattern occurs on IEEE--39.

\subsection{A sufficient gate for the complete DDE spectrum}
This standard homotopy argument prevents pole assignment from
silently ignoring the other roots. At fixed $\rho$, in a common
gauge-deflated $E=I$ state chart, let
\[
 \Delta_1(s)=\Delta_0(s)+U(s)C,\qquad
 U_i(s)=b_i^0e^{-s\tau_i^0}-b_i^1e^{-s\tau_i^1}.
\]
Both delay vectors are prescribed.
Let $\Gamma$ enclose $\Re s\ge-\sigma$, truncated by $|s|\le R$,
with $\Delta_0$ nonsingular on $\Gamma$. A sufficient condition is
\begin{equation}
 \sup_{s\in\Gamma}\|C\Delta_0(s)^{-1}U(s)\|_2<1 . \label{eq:contour}
\end{equation}
Indeed, for $0\le t\le1$ the determinant lemma and Neumann
invertibility make $\Delta_0+tUC$ nonsingular on $\Gamma$.
Its argument-principle root count is unchanged.
There are no omitted roots for $|s|>R,\Re s\ge-\sigma$ if
\[
 R>\|A_0\|_2+
 \max_{k=0,1}\sum_i\|b_i^k\|_2\|c_i\|_2e^{\sigma\tau_i^k}.
\]
Dominance of $sI$ proves this bound throughout the affine homotopy.
A zero initial root count and verified (\ref{eq:contour}) therefore
imply the full spectral margin.

The gauge must actually be removed; pinned roots on the contour
need separate deflation and are excluded here.
This sufficient condition is potentially conservative and has not
been evaluated. Sampling cannot certify its supremum.
Changing $\rho$ requires the full characteristic difference,
including $A_0,C$ and equilibrium dependence, rather than the
fixed-$\rho$ low-rank update above. Nonlinear event guards remain
separate from this spectral argument.

\paragraph{Deferred validation; no experiments in this extension.}
First check the gain-sensitivity ratio and spillover formula on
withheld simple modes. Then check finite gain transport, all
gain bounds and the complete spectrum. Only afterwards test
the frozen nonlinear events and whether a replacement predictor
improves a feasible design under identical limits. No trajectories,
optimization runs or spectral sweeps were executed here.

\section{Closing the design problem before simulation}\label{sec:closure}

\paragraph{Research contract.}
The primary question is whether the collective grid return can predict
an admissible PLL retuning as SG replacement increases. The proposed
mechanism is incompatibility between the gains demanded by different
oscillatory patterns. The falsifying test is that this incompatibility
fails to predict useful full-model designs, or disappears when unnecessary
pattern restrictions are removed. The consequence, if validated, is a
controller decision based on network interactions. Neighbor feedback is
a separately identified architecture extension, with its own latency and
gain contract. None of these statements assumes a strict advantage exists.

\subsection{Open every PLL: the plant remains the full grid}
Remove all PLL triples $(\theta,\omega_p,\xi)$ while keeping the $n$
angles as external inputs to the other device/network equations:
\[
 \mathcal A(s;\rho)z=\mathcal B(s;\rho)\theta,\qquad
 e=\mathcal C(s;\rho)z+\mathcal D(s;\rho)\theta .
\]
For $\mathcal A$ nonsingular, define
\begin{equation}
 G(s;\rho)=\mathcal D+\mathcal C\mathcal A^{-1}\mathcal B,
 \qquad D_f(s)=s^2(I+sT_f). \label{eq:allreturn}
\end{equation}
Unlike $G_i$ in (\ref{eq:return}), $G$ is independent of \emph{all}
removed PLL gains and detector delays. The current, DC, SG, lossy AC,
P/Q and equilibrium dependencies remain in $G$. This independence uses
the actual model contract: gains enter only the removed PLL equations.
Do not silently reuse it for a model with other gain-dependent channels.

For local diagonal gains $K_p,K_i$, the exact reduced characteristic is
\begin{equation}
 F(s;\rho,K)=D_f(s)-(sK_p+K_i)E_\tau(s)G(s;\rho),
 \qquad E_\tau=\diag(e^{-s\tau_i}). \label{eq:allF}
\end{equation}
Its determinant describes full-model roots wherever the eliminated
block is regular; hidden-block singularities need the original descriptor.
No approximation of the exponentials or modal diagonalization is used.

\begin{theorem}[Analytical gains for a prescribed collective pattern]
Fix a replacement $\rho$, $\lambda=\alpha+\ii\omega$ with $\omega>0$,
and a nonzero PLL-angle vector $q\in\mathbb C^n$. Assume
$\mathcal A(\lambda;\rho)$ is nonsingular and
$y_i=[G(\lambda;\rho)q]_i\ne0$ for every site.
Set
\begin{equation}
 W_i=\frac{\lambda^2(1+t_{f,i}\lambda)e^{\lambda\tau_i}q_i}{y_i},
 \qquad
 \boxed{K_{p,i}=\frac{\Im W_i}{\omega},\quad
 K_{i,i}=\Re W_i-\alpha\frac{\Im W_i}{\omega}.} \label{eq:allgain}
\end{equation}
These are the unique real local gains for which $F(\lambda)q=0$.
They are admissible precisely when they obey the declared gain bounds.
\end{theorem}
\status{EXACT\_IDENTITY}
\begin{proof}
Row $i$ of $F(\lambda)q=0$ is
$(\lambda K_{p,i}+K_{i,i})e^{-\lambda\tau_i}y_i
=\lambda^2(1+t_{f,i}\lambda)q_i$.
Its two real equations give (\ref{eq:allgain}). Reconstruct
$z=\mathcal A^{-1}\mathcal Bq$, $\omega_p=\lambda q$ and
$\xi=K_iE_\tau y/\lambda$ to recover the full eigenvector.
Real model coefficients also give the conjugate eigenpair.
\end{proof}
If $y_i=0$ and the right side is nonzero, assignment is impossible.
If both vanish, that row imposes no gain equation and must be treated
without division. Nonzero $q$ is required; its complex scale cancels
in (\ref{eq:allgain}). This result assigns an eigenpair, not dominance,
simple multiplicity, the complete spectrum, or nonlinear safety.
Conversely, any existing local-gain design with such an eigenpair
and nonzero responses satisfies this same formula. Thus allowing
the pole and pattern to vary covers this regular eigenpair chart;
holding a convenient pattern fixed is the additional restriction.

The collective dependence is explicit:
\[
 y_i=\sum_jG_{ij}(\lambda;\rho)q_j .
\]
Thus even local gains depend on the spatial amplitudes and phases of
the whole pattern. Replacing this sum by $G_{ii}q_i$ is an additional
model approximation. Offline use of $G$ does not require online
communication or make the controller distributed.

\subsection{Exact gain compatibility for several patterns}
For $m$ distinct nonreal target pairs $(\lambda_r,q_r)$, construct
$W_{ir}$ from (\ref{eq:allgain}). Let
\[
 C_\Lambda=
 \begin{bmatrix}
 \alpha_1&1\\ \omega_1&0\\
 \vdots&\vdots\\ \alpha_m&1\\ \omega_m&0
 \end{bmatrix},\qquad
 w_i=\begin{bmatrix}\Re W_{i1}\\\Im W_{i1}\\
 \vdots\\\Re W_{im}\\\Im W_{im}\end{bmatrix}.
\]
Then, with no gain bounds, simultaneous local assignment is possible
if and only if
\begin{equation}
 \boxed{(I-C_\Lambda C_\Lambda^\dagger)w_i=0
 \quad\hbox{for all }i.} \label{eq:compat}
\end{equation}
The gains are $[K_{p,i},K_{i,i}]^T=C_\Lambda^\dagger w_i$.
Since $\omega_r>0$, $C_\Lambda$ has rank two and the gains are unique;
bounded feasibility also requires membership in the gain rectangle.
\status{EXACT\_IDENTITY}
The proof is the range criterion for $C_\Lambda k_i=w_i$.
For two pairs this is four real equations for two real gains at each
site. The conflict is therefore testable, rather than inferred from
large or negative damping coefficients.

\textbf{Scope restriction.} Failure proves infeasibility for these
\emph{specified poles and patterns}. It does not prove that local
controllers cannot stabilize the same replacement with other modes.
It is not a decentralized fixed-mode theorem. Patterns and target
frequencies must be allowed to vary before claiming an architectural
stabilization limit.

\subsection{Which additional measurement can remove a conflict?}
For an admissible directed information pattern $\mathcal N_i$, including
the local signal, consider the explicit extension
\begin{align*}
 t_{f,i}\dot\omega_{p,i}
 &=\xi_i-\omega_{p,i}
   +\sum_{j\in\mathcal N_i}k_{p,ij}e_j(t-\tau_j-d_{ij}),\\
 \dot\xi_i&=\sum_{j\in\mathcal N_i}k_{i,ij}e_j(t-\tau_j-d_{ij}),
 \qquad \dot\theta_i=\omega_{p,i}.
\end{align*}
Every $d_{ij}\ge0$ is an exogenous communication delay, with $d_{ii}=0$.
Signals are calibrated detector errors with declared units, sign and
timestamps. Their zero locked values preserve the equilibrium.
Off-diagonal gains and their bounds define a new controller contract;
they must not be treated as the old twenty-gain architecture.

For fixed target $(\lambda_r,q_r)$ define
\[
 y_r=G(\lambda_r;\rho)q_r,\quad
 z_{ijr}=e^{-\lambda_r(\tau_j+d_{ij})}(y_r)_j,\quad
 b_{ir}=\lambda_r^2(1+t_{f,i}\lambda_r)(q_r)_i .
\]
The row condition is
\[
 \sum_{j\in\mathcal N_i}
  [\lambda_r z_{ijr}k_{p,ij}+z_{ijr}k_{i,ij}]=b_{ir}.
\]
Stack real and imaginary parts into the real linear system
\begin{equation}
 A_i x_i+B_i u_i=b_i. \label{eq:neighbor}
\end{equation}
$A_i$ contains the two self-gain columns; $B_i$ contains the proposed
neighbor-gain columns. A proportional column for $j$ is the stacking
of $\Re(\lambda_r z_{ijr}),\Im(\lambda_r z_{ijr})$; an integral column
stacks $\Re z_{ijr},\Im z_{ijr}$. Apply a fixed declared row scaling
before using residual norms or singular values. Exact feasibility is
unchanged by invertible row scaling.

\begin{theorem}[Compatibility and minimum added gain]
For (\ref{eq:neighbor}), let
$\Pi_i=I-A_iA_i^\dagger$, $r_i=\Pi_i b_i$ and $H_i=\Pi_i B_i$.
Assume initially that self gains are unrestricted.
There exist real gains realizing every prescribed pair at site $i$
if and only if $r_i\in\operatorname{range}H_i$.
With added-gain metric $R_i\succ0$, the unique minimum-metric neighbor
gain is
\begin{align}
 u_i^\star&=R_i^{-1}H_i^T
   (H_iR_i^{-1}H_i^T)^\dagger r_i, \label{eq:neighborgain}\\
 \kappa_{i,\min}^2&=
   r_i^T(H_iR_i^{-1}H_i^T)^\dagger r_i. \label{eq:neighborcost}
\end{align}
Self gains then solve $A_ix_i=b_i-B_iu_i^\star$.
\end{theorem}
\status{EXACT\_IDENTITY}
\begin{proof}
Project (\ref{eq:neighbor}) onto the orthogonal complement of
$\operatorname{range}A_i$. This gives $H_iu_i=r_i$, which is also
sufficient because the remaining residual belongs to
$\operatorname{range}A_i$. Whiten $u_i$ with $R_i^{1/2}$ and apply
the minimum-norm solution of a consistent linear system; this proves
both formulas. Its nullspace is orthogonal in the metric to the
displayed solution, proving minimality and uniqueness.
\end{proof}
An added budget $u_i^TR_iu_i\le\kappa_i^2$ is feasible if and only if
the range test holds and $\kappa_{i,\min}\le\kappa_i$, still with
unrestricted self gains. With bounds on either gain family, solve
the exact convex feasibility/QP problem (\ref{eq:neighbor}) with
those bounds; the displayed unconstrained minimizer is not sufficient.
Polynomial graph coefficients shared across sites require stacking
all rows into one linear problem, rather than independent row solves.

With $d_i=|\mathcal N_i|$ incoming signals, the assignment matrix has
$2d_i$ columns and $2m$ rows. Arbitrary modal demands require
$d_i\ge m$ and full row rank. One neighbor can resolve a two-pair
conflict only when its two columns supply the missing directions;
communication delays can destroy or improve this rank/conditioning.
A graph edge count alone is insufficient. These statements concern
prescribed eigenstructure, not global stabilizability.

\subsection{Why arbitrary powers of a Laplacian do not close the problem}
Let $L$ be a normalized, symmetric, connected Laplacian and $K(L)$
a polynomial graph filter. For delayed source measurements,
\begin{equation}
 [L,K(L)E_\tau(s)]=K(L)[L,E_\tau(s)]. \label{eq:polycomm}
\end{equation}
If $K(L)$ is invertible, the left side vanishes exactly when the
rightmost commutator vanishes. Its entry is
$L_{ij}(e^{-s\tau_j}-e^{-s\tau_i})$.
Commutation throughout a frequency interval therefore forces equal
delays along all edges, hence uniform delay on a connected graph.
At isolated frequencies, phase coincidence can be an exception.
Furthermore,
\begin{equation}
 \|[L,E_\tau(\ii\Omega)]\|_F^2
 =8\sum_{\{i,j\}\in\mathcal E}w_{ij}^2
       \sin^2\!\left(\frac{\Omega(\tau_i-\tau_j)}2\right).
 \label{eq:freqcomm}
\end{equation}
\status{EXACT\_IDENTITY}
Both formulas follow entrywise; each undirected edge occurs twice.
If $\sigma_{\min}(K(L))\ge\gamma>0$, the remaining commutator norm
is at least $\gamma\|[L,E_\tau]\|_F$.
This is a channel-mixing result, not a stability lower bound.
It does not apply unchanged to nonuniform diagonal gains plus
graph terms, or to general edge-delayed filters.

The physical network basis $U$ is useful for expressing $q=Ua$ and
decomposing $Gq$. Using all graph coordinates is an exact basis
change; truncating $a$ restricts the design and needs separate evidence.
The actual model need not commute with $L$, even at zero delay.
Powers $L^\ell$ encode weighted walks, not an independent physical
cycle basis. No new instability is established by such powers alone.

\subsection{A calculable replacement predictor and corrector}
Use (\ref{eq:allgain}) as an exact gain-elimination map
$K=\Phi(\rho,\lambda,q)$ on a regular chart. Normalize one nonzero
complex component of $q$ to one. For fixed delays and nonzero
$q_i,y_i,W_i$, its differential is
\begin{align}
 dW_i&=W_i\left[
 \left(\frac2\lambda+\frac{t_{f,i}}{1+t_{f,i}\lambda}+\tau_i\right)d\lambda
 +\frac{dq_i}{q_i}-\frac{dy_i}{y_i}\right],\\
 dy&=G\,dq+G_s q\,d\lambda+
             \sum_jG_{\rho_j}q\,d\rho_j,\\
 dK_{p,i}&=\frac{\Im(dW_i)-K_{p,i}\,d\omega}{\omega},\\
 dK_{i,i}&=\Re(dW_i)-\alpha\,dK_{p,i}-K_{p,i}\,d\alpha .
 \label{eq:gaindifferential}
\end{align}
Use the undivided quotient rule at a zero numerator. Derivatives of
$G$ include the algebraic closure and any equilibrium dependence.
The chain rule for another simple root is
\[
 d\mu=\sum_j\mu_{\rho_j}d\rho_j+
 \sum_i(\mu_{K_{p,i}}dK_{p,i}+\mu_{K_{i,i}}dK_{i,i}),
 \qquad
 \mu_p=-\frac{\ell^*Z_p r}{\ell^*Z_s r},
\]
where the physical-port representation must be regular; otherwise
use the full descriptor sensitivity. This retains the collective,
reactive and delay effects, including the spillover in
(\ref{eq:spillover}).

In normalized real coordinates
$\zeta=(\rho,\alpha,\omega,\Re q_{\rm free},\Im q_{\rm free})$, let
$p(\zeta)=(\rho,\Phi(\rho,\lambda,q))$.
The replacement problem on this chart is
\[
 \min_\zeta\ \sum_iP_i^0(1-\rho_i),\qquad
 g_j(p(\zeta))\le0,\quad \alpha\le-\sigma_{\rm req},
\]
with the actual gain/replacement bounds, every required physical root,
and all nonlinear event guards included in $g$. This is an exact
reparameterization on the declared chart, not a convexification of
the security constraints. A computable local subproblem is
\begin{equation}
 \max_{\delta\zeta}\ (P^0)^T\delta\rho
       -\frac{\beta}{2}\delta\zeta^T R_\zeta\delta\zeta,\quad
 g_k+J_{g,k}J_{p,k}\delta\zeta\le0,\quad
 \delta\zeta\in\mathcal T_k . \label{eq:closureQP}
\end{equation}
Here $\beta>0$, $R_\zeta\succ0$ and the convex trust region
$\mathcal T_k$ are declared in advance. It is a convex QP when
$\mathcal T_k$ is polyhedral. Include all active smooth branches of
peak/spectral constraints; a single gradient of their maximum is
not generally valid. $J_p$ uses (\ref{eq:gaindifferential}).
This is a standard SQP predictor with an analytical gain map.

The method to implement after the algebra gate is:
\begin{enumerate}
 \item Start from a fully feasible design and a simple visible oscillatory
 root, with its actual PLL pattern. Confirm reconstruction by
 (\ref{eq:allgain}); this parity test is mandatory.
 \item At a frozen $\rho$-direction, compute the predicted gain change
 from (\ref{eq:gaindifferential}). Alternatively allow $\rho$ to vary
 in a trust region and maximize $(P^0)^T d\rho$ using all active
 linearized security constraints after substituting $dK$.
 \item Permit $\omega$, the admissible real part, and the complex
 pattern to vary. Fixing them is a declared restricted predictor,
 not a fair upper bound on local-controller capability.
 \item Evaluate a finite trial through $K=\Phi(\rho,\lambda,q)$,
 rather than applying only the tangent gain update.
 This makes the chosen eigenpair exact within arithmetic/model error.
 Reject bound violations and singular or poorly conditioned charts.
 \item Check the complete required spectral region and the frozen
 nonlinear guards. If another mode becomes limiting, include it and
 use the compatibility equations with variable patterns in a local
 Newton/SQP correction. Reduce the step when predicted and actual
 constraint changes disagree.
 \item Stop at a declared local stationarity/residual criterion,
 a chart failure, or an unresolved feasibility conflict. Record which.
 Do not report a restricted-family stop as the full replacement maximum.
\end{enumerate}
For prescribed patterns, the gain solve is explicit or convex. Jointly
varying patterns and replacement makes the outer problem nonconvex.
With all $q$ coordinates free, this parameterization does not promise
dimension reduction or computational superiority. Local Newton
convergence requires a nonsingular normalized residual Jacobian;
root collisions require a subspace/nonsmooth method. No global
convergence or quantitative speedup is asserted.

No frequency, RoCoF, voltage or actuator guard is inferred merely from
pole assignment. Their original nonlinear definitions and all frozen
events remain mandatory at the later validation stage. A sufficient
full-spectrum continuation condition was given above; it has not been
evaluated. Exact local gains alone are not an exact security frontier.

\subsection{Evidence gates and literature boundary}
Pole assignment with delay, eigenstructure design, sparse feedback and
spectral optimization are established methods
\cite{ram,katewa,vanbiervliet,mao}.
Ram et al.\ treat delayed feedback in symmetric mechanical systems;
Katewa and Pasqualetti treat sparse static state-feedback assignment;
Vanbiervliet et al.\ treat nonsmooth spectral optimization for retarded
systems, including continuation of nonlinear equilibria.
These antecedents preclude claiming that the matrix solve, projection,
or iterative optimization is itself new. The narrower candidate is
the bounded PLL compatibility law under the full-grid return and its
verified consequence for replacement or latency decisions.
Recent converter-model comparisons reinforce the need to retain
hidden dynamics \cite{coffey}; supplementary GFL power-oscillation
controllers \cite{chakraborty} are a different action set from PLL-only tuning.

Before new grid simulations, close three gates:
(i) the exact gain map, compatibility and derivatives pass algebra checks;
(ii) the controller architecture, fixed delays, bounds, target-selection
rule and falsification criteria are written down;
(iii) the first validation has a numerical prediction to test, not a
desired replacement percentage. Gate (i) is mathematical only; IEEE--39
parity, all-root checks and event validation remain unexecuted here.

A first experiment should test whether the formula reconstructs known
gains and predicts a finite replacement correction on withheld modes.
Only after that succeeds should a matched-budget design comparison run.
Neighbor feedback is investigated only if the local modal conflict is
material after relaxing arbitrary pattern choices. A rigorous claim
that neighbors enable otherwise impossible secure replacement would
need a valid local-architecture upper bound below a fully feasible
distributed design; fixed-pattern incompatibility is not that bound.

\section{Claim ledger, novelty and missing gates}
\begin{center}\small
\begin{tabular}{p{.49\linewidth}p{.43\linewidth}}
\toprule
Result & Status\\\midrule
Finite rank-one impedance change and rank-two indefinite damping
under independent paths & EXACT\_IDENTITY for the full linearization\\
PLL phase tracking: delay first at cubic order & EXACT\_IDENTITY\\
Graph moment and signed-area overshoot theorem & PROVED\_REDUCED\_MODEL\\
Nonlinear pre-delay gain independence & EXACT\_IDENTITY under stated history\\
Single-site assignment and heterogeneous-delay gain transport
& EXACT\_IDENTITY; one assigned pair only\\
Quantified spillover under pole-preserving retuning
& EXACT\_IDENTITY for simple roots and stated channels\\
Gain-bounded replacement in a fixed direction
& PROVED\_REDUCED\_MODEL for first-order constraints\\
Complete DDE root-count preservation condition
& Conditional sufficient theorem; bound not evaluated\\
All-PLL eigenpair gain map and prescribed-pattern compatibility
& EXACT\_IDENTITY; hidden-block regularity required\\
Minimum added-neighbor gain for fixed modal demands
& EXACT\_IDENTITY; bounds need the stated convex solve\\
Replacement predictor with analytical gain elimination
& SUPPORTED\_LOCAL formulation; no grid run or convergence evidence\\
IEEE--39 near-zero regularity and graph-moment applicability
& BLOCKED: not established\\
Maximum replacement and globally optimal gains & BLOCKED: not obtained\\
First-ever result or superiority over literature & BLOCKED: not established\\
\bottomrule
\end{tabular}
\end{center}
Related work already treats network/PLL interactions \cite{huang},
damping--inertia--delay tradeoffs \cite{pates}, heterogeneous aggregate
frequency/nadir/RoCoF \cite{paganini}, heterogeneous inertia/control
placement \cite{pagnier}, and delayed type-II PLL design \cite{wilson}.
Their model contracts differ; these are antecedents, not reproduced
head-to-head baselines. Standard matrix identities are not novel claims.
Dominant-pole PI/PID tuning with delays and the need to verify the
other poles are also established topics \cite{zitek,zhang}.
The narrower candidate is the full-grid retuning spillover factorization
and its usefulness for gain-bounded replacement decisions.
Its bibliographic originality has not been established. The closure in
Section~\ref{sec:closure} adds a constructive assignment test and a
replacement predictor, without promoting assignment to stabilization.

The candidate thesis is: \emph{latency compensation at one PLL can
preserve a selected pole exactly while moving the other network modes
according to an explicit collective-residue law; the resulting
multimode conflicts can limit first-order SG replacement despite
remaining tuning range.} The physical damping redistribution and
graph moment provide supporting mechanisms, under their separate
assumptions. The general framework does not assume that the relevant
conflict exists on IEEE--39.

The remaining theory gates are: explicit gauge/zero-frequency
regularity; the lossy P/Q extension with the correct left/right
nullspaces; applicability of the path condition; and an actual nonlinear
storage or bounded-remainder argument for large-event security.
Taylor coefficients alone do not prove that security. No prize or
publication outcome is guaranteed.

\appendix
\section{Checks and preservation}
\path{verify_symbolic.py} checks the PLL series, a finite Schur update,
the Hermitian rank-two eigenvalue identity and the stable two-node
example using exact algebra. These checks supplement the proofs;
they are not IEEE--39 dynamic validation. No simulation, optimizer
or DDE spectral sweep is run. New files stay in
\path{experiments/theory_collective_damping_20261003/}.
Inspected source hashes are in \path{PROVENANCE.json}.
Frozen artifacts are preserved; no commit or push is made.

\begin{thebibliography}{99}
\bibitem{pates} R. Pates and E. Mallada,
\emph{Damping, Inertia, and Delay Robustness Trade-offs in Power Systems},
MTNS, 2018.
\url{https://mallada.ece.jhu.edu/pubs/2018-MTNS-PM.pdf}.
\bibitem{paganini} F. Paganini and E. Mallada,
\emph{Global performance metrics for synchronization of heterogeneously
rated power systems: The role of machine models and inertia},
arXiv:1710.07195v2, 2017.
\url{https://arxiv.org/html/1710.07195v2}.
\bibitem{pagnier} L. Pagnier and P. Jacquod,
\emph{Optimal placement of inertia and primary control: a matrix
perturbation theory approach}, arXiv:1906.06922, 2019.
\url{https://arxiv.org/abs/1906.06922}.
\bibitem{huang} L. Huang, H. Xin, W. Dong, and F. D\"orfler,
\emph{Impacts of Grid Structure on PLL-Synchronization Stability of
Converter-Integrated Power Systems}, arXiv:1903.05489, 2019.
\url{https://arxiv.org/abs/1903.05489}.
\bibitem{wilson} J. Wilson, A. Nelson, and B. Farhang-Boroujeny,
\emph{Parameter Derivation of Type-2 Discrete-Time Phase-Locked Loops
Containing Feedback Delays}, IEEE TCAS-II, 56(12):886--890, 2009.
\url{https://doi.org/10.1109/TCSII.2009.2034197}.
\bibitem{zitek} P. Z\'itek, J. Fi\v{s}er, and T. Vyhl\'idal,
\emph{Dimensional analysis approach to dominant three-pole placement
in delayed PID control loops}, Journal of Process Control,
23(8):1063--1074, 2013.
\url{https://www.sciencedirect.com/science/article/abs/pii/S0959152413001236}.
\bibitem{zhang} W. Zhang, Y. Cui, and X. Ding,
\emph{An Improved Analytical Tuning Rule of a Robust PID Controller
for Integrating Systems with Time Delay Based on the Multiple
Dominant Pole-Placement Method}, Symmetry, 12(9):1449, 2020.
\url{https://doi.org/10.3390/sym12091449}.
\bibitem{ram} Y. M. Ram, J. E. Mottershead, and M. G. Tehrani,
\emph{Partial pole placement with time delay in structures using the
receptance and the system matrices}, Linear Algebra and Its Applications,
434(7):1689--1696, 2011.
\url{https://doi.org/10.1016/j.laa.2010.07.014}.
\bibitem{katewa} V. Katewa and F. Pasqualetti,
\emph{Minimum-Gain Pole Placement With Sparse Static Feedback},
IEEE Transactions on Automatic Control, 66(8):3445--3459, 2021.
\url{https://doi.org/10.1109/TAC.2020.3018615}.
\bibitem{vanbiervliet} J. Vanbiervliet, K. Verheyden, W. Michiels,
and S. Vandewalle, \emph{A nonsmooth optimisation approach for the
stabilisation of time-delay systems}, ESAIM: COCV, 14(3):478--493, 2008.
\url{https://doi.org/10.1051/cocv:2007060}.
\bibitem{mao} X. Mao and H. Dai,
\emph{Partial eigenvalue assignment with time delay robustness},
Numerical Algebra, Control and Optimization, 3(2):207--221, 2013.
\url{https://doi.org/10.3934/naco.2013.3.207}.
\bibitem{coffey} S. Coffey, S. Harrison, A. Egea-Alvarez, and L. Xu,
\emph{Assessing the Capability and Limitations of Small Signal
Simplified Grid Following Models for the Identification of
Sub-Synchronous Oscillations}, IEEE Access, 2025.
\url{https://doi.org/10.1109/ACCESS.2025.3586197}.
\bibitem{chakraborty} T. Chakraborty, A. Pal, and S. Maleki,
\emph{Damping LFOs: Grid Following with Power Oscillation Damping
vs. Grid Forming vs. PSS}, arXiv:2505.24204, 2025.
\url{https://arxiv.org/abs/2505.24204}.
\end{thebibliography}
\end{document}
